\chapter*{Abstract}
Email APIs have become the component a modern application uses to send
email. Providers such
as SendGrid, Resend, Mailgun, and Amazon SES offer services that
enable transactional and marketing email delivery through programmatic
interfaces. However, choosing the most suitable API for a given project
is a poorly informed decision. How many requests per second an
endpoint admits, how many recipients fit into a single call, or
what happens once a quota is exceeded are figures that rarely appear on
the pricing page. They are spread across developer documentation,
support articles, and acceptable-use policies, when they appear at all.
The comparisons that are easy to find are written by the providers
themselves.

This situation makes objective and systematic comparison between APIs
difficult. Take a team that needs to send 100,000 emails a month in bursts of
1,000 requests per minute, with webhook support and batch sending. It
has no way of finding out which provider meets all of that at once. Standards such as OpenAPI exist for the
functional description of REST APIs. An API's consumption conditions, by
contrast, lack an established equivalent that captures them in a
structured, comparable, and machine-readable manner.

This work is framed within \textit{Pricing Intelligence}, the research
line that aims to turn the information contained in pricing plans into
knowledge that supports decision-making. Within that line, previous work
by the group has conceptualised the contractual capacity of APIs through
the \textit{API Capacity Model} (ACM) and its ACM2YAML serialisation.
From it comes the \textit{Datasheet}: a machine-oriented document that
brings together in a single representation the information that is
scattered. The tools of the ecosystem operate on it: PRIME4API as the
capacity analysis engine, HARVEY4API as the query interface, and SPHERE
as the publication hub. These artefacts, however, have so far been validated
against three pricings of two services \citep{gavira2026tfm}, and their
operational version captures two of the four dimensions of the metamodel.

The goal of this Master's Thesis is to take that instrumentation to a
real domain and at scale. To this end, a catalogue of Datasheets is built
for a representative set of email APIs, rigorously documented from each
provider's official technical documentation and with explicit source
traceability. Confronting the model with the domain also makes it
possible to identify and correct its expressiveness gaps.

The contribution of this work is fourfold. First, it delivers a curated,
versioned catalogue of Datasheets for the email domain, prepared as a
contribution to SPHERE and reusable by the rest of the ecosystem. Second, it proposes an
evolution of the ACM2YAML schema justified case by case: every extension
responds to a real-world scenario that the previous version was unable to
represent. Third, it characterises how usage limits and rates interact
within the domain: when the rate keeps the quota from ever being
exhausted, when the quota runs out before the rate is even felt, and what
that costs the consumer. The relationship was already defined in the
model, but not characterised on a real domain. It is analysed through the
geometric interpretation of capacity and the analysis operations provided by PRIME4API. Fourth, it establishes a
systematic comparison method that answers decision questions such as
\textit{``which API is most suitable if my requirements are X?''}. The method weighs functional and non-functional
requirements in the same query.

The brief set three sources for selecting the APIs. They were those
analysed in a previous master's thesis from the same programme, those already modelled
within the ecosystem, and a documented query to a language model. The
first could not be used, because that thesis arrived after the selection
had been made. The other two, plus the pilot of the method, yield the
eight APIs of the catalogue.

The results show that formally structuring knowledge about email APIs
makes visible differences between providers that conventional comparisons
do not reveal. They also show that a significant part of these services'
contractual capacity is never published: this marks the limit of what can
be known from documentation and leaves the rest to empirical measurement.
Both the instrument and the method are expected to extend to API domains
beyond email, although that extensibility remains untested.

\vspace{.5cm}
\noindent\textbf{Keywords:} Pricing Intelligence, email APIs, API datasheets, API capacity, rate limits, usage quotas, systematic API comparison, ACM2YAML, non-functional properties, service selection
